\documentclass{IEEEcsmag}

\usepackage[colorlinks,urlcolor=blue,linkcolor=blue,citecolor=blue]{hyperref}
\expandafter\def\expandafter\UrlBreaks\expandafter{\UrlBreaks\do\/\do\*\do\-\do\~\do\'\do\''\do\-}
\usepackage{upmath,color}
\usepackage{lipsum}
\usepackage{changebar}

%%\usepackage{draftwatermark}
%%\SetWatermarkText{DRAFT}


\jvol{XX}
\jnum{XX}
\paper{8}
\jmonth{Month}
\jname{\textit{IEEE Annals of the History of Computing}}
\jtitle{\textit{IEEE Annals of the History of Computing}}
\pubyear{2026}

\newtheorem{theorem}{Theorem}
\newtheorem{lemma}{Lemma}
\newcommand\blfootnote[1]{%
  \begingroup
  \renewcommand\thefootnote{}\footnote{#1}%
  \addtocounter{footnote}{-1}%
  \endgroup
}

\setcounter{secnumdepth}{0}

\begin{document}

\sptitle{ARTICLE}

\title{The First Fifteen Years of Computing in Aotearoa New Zealand}

\author{Brian E. Carpenter}
\affil{University of Auckland, New Zealand}

\author{Sathiamoorthy Manoharan}
\affil{University of Auckland, New Zealand}

\author{Janet Toland}
\affil{Victoria University of Wellington, New Zealand}


\markboth{ARTICLE}{ARTICLE}


\begin{abstract}
\cbstart
This article reviews and discusses the first fifteen years of electronic computing in New Zealand.
Its goal is to consider whether the country's specific nature affected the deployment of computing compared to other countries, and conversely whether computing uniquely affected post-colonial New Zealand. We approach these questions by documenting  
the growth in computers and staff, types of usage including service bureaus, the professionalization of the industry, educational and academic developments, and examples of local ingenuity, during 1960-1975. We discuss the economic and social impact of computing during a period of major economic and social change for the country, as it relinquished its dependency on the United Kingdom. We then compare New Zealand with various other countries to establish whether New Zealand had any claim to exceptionalism during the original dissemination of modern computing.
\cbend
\end{abstract}


\maketitle


\chapteri{M}odern computers arrived in Aotearoa New Zealand in 1960 \cite{FirstCinNZ}. 
New Zealand was then a relatively isolated and small economy, with a growing population
of about 2.5 million, so it is feasible to track the dissemination and socio-economic influence
of computing for the following period; this study runs from 1960 to about 1975. Two events
in the mid-1970s signaled the end of the economic and social era that began after World War II.
Firstly, Britain joined the European Common Market in 1973, with major economic consequences
for New Zealand.
\cbstart
This coincided with the first global oil price shock, ending an era of comfortable post-colonial prosperity.
Secondly, New Zealand switched on a national Police Computer in 1976,
ending a period during which computers were largely invisible to the general
population and their social impact was minimal. At about the same time, the computer
market changed its nature due to the increasing prevalence of minicomputers.

The main questions we wish to address are whether, during our period of study,
the specific nature of New Zealand’s economy and society significantly affected
\cbend  %%this bizarreness works around an apparent LaTeX bug
the deployment of computing technology compared to other countries,
\cbstart  %%this bizarreness works around an apparent LaTeX bug
and conversely whether this deployment
uniquely affected New Zealand’s post-colonial society.
New Zealand often claims some degree of cultural exceptionalism, e.g. \cite{Frykberg-2025}; we wish to question whether this was justified during the early years of modern computing. As a basis for tackling these questions, we start by presenting factual material on various aspects, relying on contemporary sources, existing historical studies, and the memories of some participants:
\cbend

\begin{itemize}
\item growth in numbers of computers, vendors and employment
\item key areas  and types of usage, including service bureaus
\item professionalization
\item training, teaching and academia
\item local ingenuity
\end{itemize}

\cbstart
In the years we are studying, New Zealand’s economy was booming, at the same time that its transformation from a Dominion within the British Empire to an independent nation was changing its society as a whole. Thus, we will
consider the economic importance of computing, and its impact on society, during this time of change. Finally, we compare New Zealand with other developed countries, to consider whether the country exhibited any unique properties. After fifteen years of computing, was New Zealand lagging or leading by comparison with its peers?
\cbend
\vspace*{-8pt}
\section{GROWTH IN NUMBERS AND EMPLOYMENT}
Like other developed countries, early forms of information technology had appeared in New Zealand during the first half of the 20th century, but were of moderate importance in its largely agricultural economy \cite{pioneers}.
By 1960, the country had a rather centrally managed post-colonial economy. Indeed, its currency was tied to the pound sterling until as late as 1985, so the country remained somewhat bound to the British economy. The details are complex \cite{Sullivan2013}, but the result was that during the period of our study the Treasury was always concerned about the foreign exchange impact of computer imports. This necessitated a
\cbstart
system of import licensing that had a major impact on the early years of computing, since even the richest companies could not import computers at will -- the market was significantly under-served.
\cbend 
For example, in the financial year 1966/67 requests to import computers totalled over NZ£2,000,000 in value but only NZ£400,410 was available \cite{Treasury-1966}. Priority was given to companies which brought in the most foreign exchange. This was a constant background for trends in computing, especially since the Treasury initially favored computing service bureaus to ensure maximum usage of what they considered a scarce resource. Of course, many companies much preferred to have their own systems.

Table~\ref{TotalComp} shows the estimated number of electronic digital computers installed in New Zealand over the years in question.

\begin{table}[h]
\begin{center}
\begin{tabular}{ |c|c|c|c| } 
 \hline
 Year & Computers & Source & Notes\\ 
\hline
1960 & 2 & \cite{FirstCinNZ}&\\ 
1965 & 70 & \cite{Beardon} & \cite{HoneHeke} gives 45\\ 
1966 & 81 & \cite{Beardon}&\\
1967 & 100 & \cite{ScottSurvey1967}&\\
1968 & 120 & \cite{Beardon} & \cite{Yearbook75} gives 87\\
1969 & 140 & \cite{Beardon}&\\
1971 & 180 & \cite{Beardon}&\\
1972 & 200 & \cite{Yearbook75}&\\
1974 & 280 & \cite{Beardon}&\\
1976 & 400 & \cite{Beardon}&\\
 \hline
\end{tabular}
\vspace{1ex}
\caption{\label{TotalComp}Estimated number of electronic digital computers installed in New Zealand.}
\end{center}
\end{table}
A snapshot of installations in Auckland alone lists 37 sites at the end of 1969, covering businesses of many sorts, local government, a university, and even the Naval Research Laboratory \cite{Auburn1971}. At least seven of the sites appear from their names to be service bureaus. The author noted in passing that there was also a large Government investment in computers in Wellington.

One example was the ICT 1900 machine installed in 1966 by Winstone Ltd at 69 Queen St, Auckland. Andy and Tom Bell \cite{Bell-private} recall that their British immigrant father William (Bill) Bell, an ICT and former Powers-Samas technician, supported this machine, and the young brothers were allowed to visit the machine room (Figure~\ref{Winstone}). Winstone was later taken over by the Fletcher Building company.

\begin{figure}
\centerline{\includegraphics[width=18.5pc]{Winstone1966.jpg}}
\caption{\label{Winstone}Winstone computer room, 1966. Courtesy Fletcher Archives, NZ.}\vspace*{-5pt}
\end{figure}

Although there are some discrepancies in the available data, the trend is one of rapid growth, from one computer for every 1.25 million people to one for every 10,000. Similarly, the number of programmers in public services rose from a handful (maybe 4) in 1960 to 115 (84 men and
31 women) by 1974 \cite{Yearbook75}. Overall, some 4000 data processing staff were then employed at  220 sites -- 39\% in Wellington, 31\% in Auckland, 9\% elsewhere in the North Island, and 21\% in the South Island \cite{Beardon}. Pay was attractive: in 1967, an EDP manager’s average annual salary was NZ\$4580 \cite{ScottSurvey1967}, when the national average for company directors and managers was NZ\$4220 \cite{Yearbook70}. Programmers and systems analysts earned anywhere between NZ\$2000 and NZ\$4000. 
\cbstart
It is impossible to tell how many of these staff were recruited locally and how many were immigrants, especially those from the UK for whom visa requirements were non-existent until 1974, so no statistics on their skills are available. Nevertheless, it is noticeable that among our contemporary informants for this article and for \cite{pioneers}, approximately a quarter were born in the UK.
\cbend

Even if the public services had 27\% of female staff, it is not clear that the same applied elsewhere. For example, a 1966 staff photograph of Fletcher Building’s internal computer bureau shows 14 men and only three women (Figure~\ref{Staff1966}). Whenever staff appear in computer room photographs, the women are almost invariably seated at a keyboard and the men are standing in discussion. A 1967 survey \cite{ScottSurvey1967} shows systems design and programming as 90\% male occupations, operations as 50\% male, and data preparation as 97\% female. 

There is little to indicate significant Māori involvement in computing during our period of study.
Robyn Kāmira \cite{Kamira} shows that such involvement would only start in the 1980s, when a first small cohort
of Māori information technology specialists graduated.

\begin{figure}
\centerline{\includegraphics[width=18.5pc]{FletcherStaff1966.jpg}}
\caption{\label{Staff1966}Fletcher Computing Bureau staff, 1966. Courtesy Fletcher Archives, NZ.}\vspace*{-5pt}
\end{figure}

New Zealand had a long history of preferring to import British products of all kinds, including electrical and electronic technology.
\cbstart
For example, in 1960 imports from the UK were 44\% of the total, compared with 10\% from the USA \cite{Yearbook62}. This imbalance did not long survive in the computer market.
\cbend
IBM was of course a major player; not only did they secure the Treasury as their first customer for an IBM 650 in 1960 \cite{FirstCinNZ}, but they also delivered a heavily discounted IBM 1620 to the University of Canterbury in 1962, the first machine in academia \cite{Dale-Canty2}. This was in line with IBM’s world-wide practice of educational discounts, a powerful marketing tool that exposed many students to the IBM brand. As time went on, IBM sold or leased numerous machines -- more IBM 1620s, at least one IBM 1401, and several IBM 1130s, until the game-changing IBM 360 series came along in 1964.

The main British competitor in 1960 was ICT, but a potential customer would certainly have been confused by the fragmented state of the British computer industry at that time \cite{Hendry-1989}.
Nevertheless, ICT sold one ICT 1201, at least six ICT 1301s, and numerous machines from the ICT 1900 series.
It was not until the various British companies were unified as ICL in 1968 that things settled down, but by then IBM had the lion’s share of the New Zealand market and ICL was falling behind. In 1974, the new ICL 2900 series had to compete with the well established and upwards compatible IBM 360 series.

Other companies that were active in the early days were Burroughs (with its E2000, a computer specialized for accounting), and NCR (Series 500, also specialized for ledgers). Another factor was that in 1967, New Zealand moved from sterling pounds, shillings and pence to decimal dollars. This conversion accelerated the adoption of computer equipment, with the Decimal Currency Board providing NZ£91,500 in grants for organizations to switch from programmable calculators to computers \cite{Treasury-1966}. This subtly affected the market, since British vendors lost the advantage of their expertise in supporting pounds, shillings and pence. Even Olivetti entered the New Zealand market by about 1970 with their compact and elegantly designed Programma series, aimed at smaller businesses \cite{Taupo-Olivetti-1970}. 

We have collected available information about individual computers present in New Zealand during the period under study from a
\cbstart
variety of sources, using formal publications, personal communications, photographic archives, sundry web pages, press reports and at least one student newspaper, but at best these cover a small sample of the market.
The details are available on-line as a spreadsheet \cite{SpreadsheetNZ}.
At the time of writing, we have identified more than 80 machines (including programmable calculators) with arrival dates up to 1975, but the estimated total for that year exceeded 280. Many of the earlier machines had of course been scrapped by then, so our sample certainly covers less than 25\% of the market. The distribution of this sample across manufacturers is shown in Table~\ref{TotalManuf}. The total for the British companies that ended up as part of ICL is 29, similar to the total for IBM. It is clear that American suppliers had become dominant, largely overcoming New Zealand’s traditional leaning towards the UK.
\cbend
\begin{table}[h]
\begin{center}
\begin{tabular}{ |c|c| } 
 \hline
Manufacturer & Count\\ 
\hline
Burroughs & 12\\
DEC & 10\\
HP & 2\\
IBM & 28\\
NCR & 2\\
\hline
EELM & 3\\
Elliot & 1\\
ICL & 8\\
ICT & 17\\ 
Olivetti & 2\\
\hline
\end{tabular}
\vspace{1ex}
\caption{\label{TotalManuf}Sample of market distribution (US and European vendors).}
\end{center}
\end{table}

Throughout the period under study, most computers remained large and comparatively heavy, even as they moved from “first generation” (vacuum tubes) through “second generation” (discrete transistors) and “third generation” (integrated circuits) to the “fourth generation” (very large scale integration). As the technology got smaller, mainframe computers became more powerful rather than more compact; space and air conditioning requirements did not substantially change. For example, the third computer installed in 1969 by Motor Specialties in Auckland, an ICL 1902A, needed to be lifted to the top floor by crane, like its predecessors, an ICT 1201 and an ICT 1301 (Figure~\ref{MS1902A}). Only after 1969 when minicomputers such as the HP 2100 and DEC PDP-11 series entered the market did this change, so that some computers could be installed almost as easily as domestic appliances.

Originally, much equipment was rented; in 1967, 70\% of “owners” actually rented their equipment, and a further 15\% rented part of it \cite{ScottSurvey1967}. By the late 1960s, the growth in installations led to a minor building boom of dedicated computer centers, which tended to be stark concrete buildings with few windows (Figure~\ref{TAB-ChCh}). A rare exception was Ceramic House in Auckland (Figure~\ref{Ceramic}), specifically designed to accomodate an ICT 1901 delivered in 1969. This building was designated as a heritage site in 2024.

\begin{figure}
\centerline{\includegraphics[width=18.5pc]{MotorSpec1902A-1969.jpg}}
\caption{\label{MS1902A}ICL 1902A delivery to Motor Specialties in 1969. Courtesy Fletcher Archives, NZ.}\vspace*{-5pt}
\end{figure}
\begin{figure}
\centerline{\includegraphics[width=18.5pc]{TAB-Chch-1970.jpg}}
\caption{\label{TAB-ChCh}The TAB Computer Centre, Christchurch, 1970. Courtesy Fletcher Archives, NZ.}\vspace*{-5pt}
\end{figure}
\begin{figure}
\centerline{\includegraphics[width=18.5pc]{CeramicHouse1969.png}}
\caption{\label{Ceramic}Ceramic House, 1969. Courtesy Auckland Libraries Heritage Collections.}\vspace*{-5pt}
\end{figure}
This section has described how computers were imported, housed, installed and staffed during the period of study. We will now consider what they were used for.
\vspace*{-8pt}
\section{HOW COMPUTERS WERE USED}
In 1960, the main expected workload was the payrolls for teachers and public servants, but expansion of the fields of use came very quickly. By the end of 1966 \cite{Treasury-1966}, Government usage included payroll, statistics, science, NZ Railways, Inland Revenue, Post Office, Social Security, Reserve Bank, Bank of New Zealand, New Zealand Broadcasting Corporation, Apple \& Pear Board, and the Dairy Board. In a surviving list of import license applications between August 1965 and May 1966 \cite{Treasury-1966}, there are three applications from local government, three from universities, but 30 from business and primary and secondary industries. The New Zealand economy took to computers rather quickly, held back mainly by import license restrictions as mentioned above.

Applications of importance for business users were listed in 1967 as invoicing, payroll, general ledger, stock control, expense analysis and sales statistics \cite{ScottSurvey1967}. All the same, mathematical, scientific and project planning applications such as PERT (Program Evaluation and Review Technique) were increasingly used. Users saw the advantages of computerization as speed, accuracy, improved reporting, improved organizational control and staff savings. Quite surprisingly, in 1967 most computers outside of academia were underutilized, with 250 operating hours per month on average, and single-shift operation still being typical. Only 68\% of installations charged internal users for usage: “one of the chief reasons why installations can be uneconomic” \cite{ScottSurvey1967}.

Between 1967 and 1970, the government carried out three surveys on electronic data processing in the public sector, which indicated an increasing shift from accounting and management applications to more scientific and technical areas \cite{Lythgoe-1970}. 
In 1973, the Department of Statistics provided the following summary \cite{Yearbook73}:

\begin{quotation}
\noindent{The demand for computers has come from Government departments, local authorities, universities, primary producer boards, private firms in industries such as printing, forestry, insurance, oil, food processing, electrical equipment manufacturing, building and construction, clothing, engineering, airways, banking, retailing, motor assembly, paint manufacturing, and stock and station agents.}
\end{quotation}

\noindent{Because of government reluctance to issue import licenses, numerous service bureaus arose to allow resource-sharing. By 1967, when the total number of computers was about 100, at least 166 companies were using bureau services operated by computer vendors, and another 78 were using independent bureaus \cite{ScottSurvey1967}. Table~\ref{Bureaux} lists some example bureaus and their approximates dates of foundation, mainly from contemporary press reports and advertisements, but the list is far from complete.}

\begin{table}[h]
\begin{center}
\begin{tabular}{ |p{0.07\linewidth}|p{0.23\linewidth}|p{0.23\linewidth}|p{0.23\linewidth}<{\raggedright}| } 
 \hline
 Year & Name & Location & Notes\\ 
\hline
1961 & Electronic Data Processing Ltd & Auckland & Closed 1963\\\hline
1963 & Electronic Data Systems Ltd & Wellington, later Auckland \& Christchurch & 50\% EELM, 50\% NZ Truth (until 1968)\\\hline 
1965 & Computer Bureau Ltd & Christchurch, later nationwide & Renamed Datacom in 1984\\\hline
1966 & Allied Computer Processors & Dunedin & \\\hline
1967 & Databank Systems Ltd & multiple & Banking consortium\\\hline
1968 & Electronic Data Processing NZ Ltd& Wellington&\\\hline
1968 & NCR Bureau & Wellington &\\\hline
1968 & Burroughs Bureau & Wellington &\\\hline
1968 & ICL Centre& Auckland&\\\hline
1969 & Computer Activities& Auckland&\\\hline
1969 & Computer Systems& Auckland&\\\hline
1969 & Data Enterprises Ltd& Tauranga& Paul and David Walker\\\hline
\end{tabular}
\vspace{1ex}
\caption{\label{Bureaux}Example computer bureaus in New Zealand.}
\end{center}
\end{table}
The first bureau listed remains something of an enigma. According to ICT’s records in the UK, it bought an ICT 1301 for “service work” \cite{ccs-t2x1}, but we have found no local records to substantiate this, beyond the formal records of the formation and dissolution of the company. Another phantom bureau hosted the government Education Department’s second computer, also an ICT 1301, leased from ICT at the end of 1963 and installed in ICT premises in Wellington. When the Minister of Education pressed a button to inaugurate it, it was said to be “New Zealand’s first computer bureau” \cite{NatEd1963}, although all indications are that it was dedicated to the Education Department. It seems that obtaining an import license for a shared bureau computer was significantly easier than for a single-user machine. The other bureaus were real enough, and we now present a few brief case studies.

Electronic Data Systems Ltd (EDS), founded in 1963, was jointly owned by the British company English Electric Leo Marconi (EELM) and the NZ Truth newspaper. It had no connection with the later American company of the same name. EELM, one of the numerous firms that ended up as part of ICL, presumably wanted to enter the New Zealand market. Apart from its earlier DEUCE and LEO computers, it made a family of KDF machines, including the KDF6, of which at least 12 were sold \cite{Hendry-1989}. The KDF6 sales records are scanty \cite{ccs-n2x1}, but  EDS had such machines in Wellington and Auckland \cite{Harsant-private}. Indeed we have photographic evidence of the machine in Auckland (Figure~\ref{kdf6}).

\begin{figure}
\centerline{\includegraphics[width=18.5pc]{KDF6-2b2s.png}}
\caption{\label{kdf6}EDS KDF6, Princes St, Auckland in about 1970. Courtesy Neil Harsant.}\vspace*{-5pt}
\end{figure}
NZ Truth, described by the National Library web site as “one of the country’s most colourful, controversial and popular newspapers” ran from 1905 until 2007. It is unclear why it decided to enter the computer bureau business in 1963. Perhaps it was simply the easiest way to get an import license for their first computer. In any case, EDS prospered, and opened a third branch in Christchurch in 1970, later switching to Burroughs and ICL equipment \cite{Press-EDS}. By 1972, the Christchurch and Wellington offices were linked by a remote job entry system. EDS remained in business in one form or another until 1994, decades after EELM no longer existed and NZ Truth had sold its stake.

Another success story started in 1965. G. Bernard Battersby (1925-1997) was the Dean of the Faculty of Commerce at the University of Canterbury, and a government advisor. Yet he joined with Paul M. Hargreaves (1938-2014) to found Computer Bureau Ltd (CBL) in Christchurch \cite{Pioneers85}. Its first computer was an ICT 1902 delivered in 1966; with a staff of seven, CBL was unprofitable until 1969, when a branch was opened in Wellington. In early 1970, CBL had some 50 customers, and the Christchurch computer was working 140 hours per week. Thus an ICL 1902A was delivered to Christchurch, the orginal 1902 was moved to Wellington, and an additional ICL 1902 was ordered for another new branch in Hamilton (Figure~\ref{CB-delivery-1971}). In Auckland, CBL had taken over Fletcher Building’s in-house computer bureau, also ICL-based. All training of systems analysts and programmers took place in Christchurch, with total staff expected to reach 70 during 1970 \cite{Press-CBL-1}. By 1974, CBL needed to order three ICL 1902T models for Auckland, Wellington and Christchurch, with a nationwide total of six computer centers and 125 staff. Furthermore, they told the newspapers \cite{Press-CBL-2}:

\begin{quotation}
\noindent{[CBL] is unique in that many of its major users are also shareholders, a number of whom previously operated their own inhouse computers. These include Fletcher Holdings, Ltd, Wilson and Horton, Ltd, the New Zealand Co-op Dairy Company, Bing Harris Sargood, Ltd, United Dominion Corporation, Ltd, and in Christchurch, the Canterbury Building Society, the S.I.M.U. Mutual Insurance Association, and the City Council.
}
\end{quotation}

\noindent{Clearly, Computer Bureau Ltd had found a winning formula. In 2025, they still prospered under the name of Datacom, claimed to be the largest technology company in New Zealand, employing at least 6000 people at offices across the Asia-Pacific region (Figure~\ref{Datacom25}).}

\begin{figure}
\centerline{\includegraphics[width=18.5pc]{CBWaikato-delivery-1971.jpg}}
\caption{\label{CB-delivery-1971}Computer Bureau delivery in Hamilton, 1971. Courtesy Fletcher Archives, NZ.}\vspace*{-5pt}
\end{figure}

\begin{figure}
\centerline{\includegraphics[width=18.5pc]{DatacomAuckland20251102.jpg}}
\caption{\label{Datacom25}Datacom, Auckland, 2025. Courtesy Brian Carpenter.}\vspace*{-5pt}
\end{figure}

\cbstart
An important but highly specialized bureau was Databank Systems Limited, founded in 1967 \cite{Toland-Databank, Archibald}. It was set up by a consortium of major banks to automate bank services and to provide a management information system for the trading banks. It established six computing centers across New Zealand, all using IBM 360 series equipment, and by late 1969, every branch of the five trading banks in the country was using Databank. By 1977, Databank had over one thousand staff; it survived until 1994, when it was sold to the Texas-based company EDS.
\cbend

Thus, as computer usage spread throughout business and government, import license pressure encouraged resource sharing and successful bureau operations. This pressure was not to be relieved until major economic reforms in the 1980s, after our study period.
\vspace*{-8pt}
\section{PROFESSIONALIZATION AND ACADEMIA}
\subsection{The New Zealand Computer Society}
\vspace*{5pt}
Professionalization started early. The New Zealand Computer Society (NZCS) was founded in the same year that the first modern computers arrived in the country \cite{Moon2010}.
Originally known as The Data Processing and Computer Society, NZCS was registered on 6 October 1960. Initial members were enthusiasts from the academic and business communities, who were aware of the impact computing was having overseas. In the early days academics had the stronger voice and there were restrictions on the numbers of members from the “machine houses” due to an unfounded fear that they would use the Society for promotion rather than learning.

The Computer Society had two complementary roles. The first was to look after the interests of its members; the second involved meeting a broader obligation to society in general. The Society wanted to be “the independent and authoritative source of information on matters to do with computing and its social impact” . This was attractive since self-regulation could be a way of forestalling legislative constraints.
Wellington (the capital) was the first branch, followed by Canterbury, Otago and Auckland. Initial activity was very much around the branch level. By 1965, there were over 200 members and the possibilty of holding national conferences and publishing a journal was first raised. 
In 1968, the Society organised its first national conference, which was so well attended that it became a regular occurrence. As well as the presentation of papers, this conference was an important opportunity for computer vendors to present their wares. Around 500 people, on average, attended each of the conferences. An overview of the first five conferences is provided in Table~\ref{NZCSconf}. As well as local keynote speakers, well known figures in the computing world from the USA, the UK and Australia also spoke at NZCS conferences. For example, at the 1970 conference, Richard Hamming from Bell Telephones, USA, spoke on “Computers and Computing in the 1970s”, in 1972 Paul Armer from Stanford University covered “Computers in Society” and in 1976, ACM President Jean Sammet spoke on “Responsibilities of Computing Professionals”. These keynotes were a major draw at conferences and widely reported on in the local media.
\cbstart
Thus, the conferences contributed to increasing public awareness that computer usage was relevant to society as a whole, and was not just  a matter for specialists. For example, in 1970 Hamming was quoted in the media as saying "People who could not adapt themselves to the change which would come about by the year 2000, when much of life would be determined by machines, would have a difficult time.” \cite{Press-Hamming} In 1972, the NZCS was reported to be against a national personal-identifier system \cite{Press-Nat-ID}, and its work on ethics and privacy generated various news reports.
\cbend

\begin{table}[h]
\begin{center}
\begin{tabular}{|p{0.07\linewidth}|p{0.20\linewidth}|p{0.26\linewidth}|p{0.12\linewidth}|p{0.05\linewidth}<{\raggedright\arraybackslash}|} 
 \hline
 Year & Location & Theme  & Papers & Ref.\\ 
\hline
1968 & Auckland & No theme & 24 & \cite{NZCS-1}\\\hline
1970 & Dunedin& No theme &  35 & \cite{NZCS-2}\\\hline
1972 & Palmerston North & Computers~in the Community &  32 & \cite{NZCS-3}\\\hline
1974 & Christchurch & Practical Computing &  20 & \cite{NZCS-4}\\\hline
1976 & Hamilton & Responsibility &  19 & \cite{NZCS-5}\\\hline
\end{tabular}
\vspace{1ex}
\caption{\label{NZCSconf}Early NZCS conferences.}
\end{center}
\end{table}
By the 1970s, the NZCS was accepted as the voice of computer practitioners in the country by both government and the public.  That recognition came with responsibilities and the Society started to examine social issues around the use of computing, as well as the quality of computing education on offer.  A Certificate in Data Processing was set up, and the Society also developed a Code of Ethics.
The foundation President of NZCS (1960-1967), Professor Gordon Oed,  is an excellent example of the links between early computing, academia and the business world. He was a senior lecturer in accounting at Victoria University of Wellington, and formerly chair of a Wellington firm of accountants \cite{Toland-2022}. Dr Bernard Battersby, previously mentioned as the co-founder of Datacom (formerly Computer Bureau Ltd.), was also NZCS President from 1967-1969. He noted that Datacom had to wait until his term as President was over before making its first profit.

Later, NZCS rebranded itself as Information Technology Professionals New Zealand (ITPNZ) \cite{Matthews}, and was eventually liquidated in 2025 due to a financial crisis.
\subsection{Training and Academia}
\vspace*{5pt}
Naturally, the need for training was an issue from a very early stage. By 1967, it was a serious concern: “It is not sufficient for the Universities to provide computer facilities and training only at the scientific level. The Technical and Management Institutes cannot be left to their own devices... the secondary schools (and before too long the primary schools) must introduce computer knowledge” \cite{ScottSurvey1967}. In practice, most training originally came from computer vendors, or people learned on the job. The first NZCS national conference in 1968 led to a committee on training requirements, which recommended computer science undergraduate courses at every university, as well as postgraduate and extramural courses. The committee “saw it as the primary role of Computer Science departments to educate system programmers, as well as future researchers” and recommended that “Commerce Departments should deal with teaching system analysts and managers” \cite{Otago-history}.

The first complete tertiary course was offered by the Auckland Technical Institute (ATI, now Auckland University of Technology) in 1967, teaching Fortran, RPG, and COBOL, along with mathematics, statistics, and accounting \cite{Bell}. Although ATI had acquired a third-hand ICT 1201 in 1966 \cite{FirstCinNZ}, student programs were processed overnight by a local company. By 1970, various technical colleges were offering courses in computing and programming, but with only 46 students enrolled nationally \cite{Yearbook73}.

As universities acquired their own computers from 1962 onwards, they naturally started offering courses on the use of computers for mathematics, statistics, science in general, and business. Often, such courses were taught by computer center staff rather than by academic faculty. Some teaching of computing in secondary schools started from 1971, when Christ’s College in Christchurch acquired a DEC PDP-8/E \cite{Stratford-06}. However, “finding computers that students could program was a challenge, and there was no national initiative around computers in schools until the end of the 20th century” \cite{Bell}. Indeed, Massey University was offering extramural computer science courses by the mid-1970s, and the most eager participants were secondary school teachers.

Among degree-granting institutions, the University of Canterbury was the first to acquire a computer, in 1962. Table~\ref{FirstUni} lists the first computers at each university. 
\begin{table}[h]
\begin{center}
\begin{tabular}{ |p{0.07\linewidth}|p{0.23\linewidth}|p{0.23\linewidth}|p{0.23\linewidth}<{\raggedright}| } 
 \hline
 Year & University & Model & Notes\\
\hline
1962 & Canterbury, Christchurch & IBM 1620 &\\\hline
1963 & Auckland & IBM 1620\textsuperscript{\textdagger} &\\\hline 
1965 & Victoria, Wellington & Elliot 503 & Shared with DSIR\\\hline
1965 & Massey, Palmerston North & IBM 1620 &\\\hline
1966 & Otago, Dunedin & IBM 360/30 &\\\hline
1969 & Waikato, Hamilton & IBM 1130\textsuperscript{\textdagger}&\\\hline
\end{tabular}
\vspace{1ex}
\caption{\label{FirstUni}First university computers.}
\end{center}
\end{table}
\blfootnote{\textsuperscript{\textdagger} This is preserved at the Bob Doran Museum of Computing.}

Various other machines were purchased, culminating in a joint purchase of five Burroughs B6700 mainframes for Auckland, Massey, Wellington, Canterbury and Otago in 1971 \cite{Good-Moon-1972}. Waikato University and Canterbury’s Lincoln College campus were provided with under-performing remote access to Auckland and Christchurch respectively. Both responded by buying their own systems, Neil Mountier at Lincoln installing both an IBM 1130 and a DEC PDP-11/40 \cite{Mountier-1990}.

The original university machines were primarily used in support of research in various disciplines; administrative use came years later. For example, in Auckland the IBM 1620, and an IBM 1130 that joined it in 1967, were much used by economists, who were “often the heaviest user,” running simulations of econometric models in 1964-70. Applications, written in Fortran, included dynamic linear models of the New Zealand economy, dynamic simultaneous equation models, and systems of linear stochastic differential equations \cite{Phillips-UoA}. The researchers were students and colleagues of Professor Rex Bergstrom, who had cut his teeth on the Cambridge EDSAC in 1952-54. Of course, mathematicians, scientists such as X-ray crystallographers, and engineers were also early users.

At Canterbury, the IBM 1620 was soon “heavily used by a wide range of university departments, frequently logging more than 500 hours a month... By July
1964, after 2 years of use, the University’s computer had been used in more than 120 projects by 14 departments,
and more than three-quarters of these projects could not have been tackled without a computer” \cite{Dale-Canty2}.

As well as research, these early machines were also used for undergraduate programming courses. However, in the 1960s, computing was mainly taught as a useful skill, not as an academic subject in its own right. The discipline of Computer Science (CS) was late to arrive in New Zealand. For example, in the UK (New Zealand’s main source of academic immigrants at that time), the Mathematical Laboratory in Cambridge had offered a post-graduate Diploma in Numerical Analysis and Automatic Computing (later renamed Diploma in Computer Science) since 1953, and the Department of Computer Science at the University of Manchester was founded in 1964, graduating its first students in 1968. CS departments also emerged in the USA during the 1960s. It was not until 1973 that three CS departments were founded in New Zealand, at the University of Canterbury (Professor John Penny),  Massey University (Professor Graham Tate) and Waikato University (Professor Darryl Smith) \cite{Bell, Waikato-50}.
CS departments at Auckland, Wellington and Otago were not established until the 1980s, outside our main study period. In Auckland at least, this reluctance was caused partly by considerations about duplication of effort, but also by assertions that Computer Science was not a “proper” science -- one view was even that computing was merely a branch of physics, according to Professor John Butcher \cite{Butcher-private}. At Otago, Brian Cox was appointed as “Lecturer in Charge of the Computer Centre” in 1964, even though the first computer, an IBM 360/30, did not arrive until 1966. Cox had used the EDSAC 2 in Cambridge as a physics graduate student, trained by Maurice Wilkes and David Barron. Although computing courses were taught from 1966, and Computer Science papers from 1972, Cox had to wait until 1984 to become the founding Professor of Computer Science \cite{Otago-history}.
\cbstart
\vspace*{-8pt}
\section{LOCAL INGENUITY}
Unlike Australia \cite{Philipson-2017}, there appear to have been no early attempts to build a complete digital computer in New Zealand. However, during the 1960s and 1970s, computing in New Zealand, especially in the universities, was characterized by adaptation and experimentation. Limited access to imported hardware, combined with a shortage of technical support and software, encouraged a culture of innovation. Systems were modified, extended, and repurposed to meet local research and teaching needs, reflecting both necessity and ingenuity. Most of these efforts are known only by anecdotal evidence, e.g. \cite{Brownlee-private, Thomas-private}. The technical details are of limited historical interest, but the following summarized examples illustrate how the famed “Kiwi ingenuity“ \cite{ingenuity}
appeared in the early days of computing.

\begin{itemize}
\item At the University of Auckland, a home-made attachment to the IBM 1620 woke up overnight users if the machine halted during a job.
\item Also at Auckland, a home-made real time clock and a floating point unit were attached to the IBM 1130, and
the 1130’s operating system was locally patched to improve performance.
\item In 1966, Bruce Moon, then at the University of Canterbury, published the first locally written programming textbook \cite{MoonProgramming}.
\item In 1968, Progeni (originally Systems \& Programs Ltd), widely regarded as New Zealand’s first software house, was founded \cite{Harpham2010}.
\item In 1975, Massey University  used a DEC PDP-11/10 as a terminal concentrator connecting low-cost VT05 terminals to its Burroughs B6700 mainframe.
\item In 1976 the Massey University Simple Instructional Computer (MUSIC) and its assembler language MUSICAL \cite{MUSIC} were simulated on the B6700.
\item Also in the 1970s, SMALL (Small Machine Algol-Like Language) was developed on the University of Auckland’s B6700. 
\item In 1975-76,  Massey University and Victoria University of Wellington established Kiwinet, a pilot project for sharing resources between their B6700 systems~\cite{ConnectingClouds}. (Full Internet connectivity would not arrive until 1989.)
\end{itemize}

Thus, the early decades of computing in New Zealand were marked by adaptation and experimentation; limited resources compelled technical creativity.
\vspace*{-8pt}
\section{AN ECONOMIC NECESSITY}
\cbend
In 1960, the New Zealand economy was very tightly coupled to the UK, despite being literally a world away.
The New Zealand currency was held at par with the British pound, and the UK accounted for 53\% of exports and 44\% of imports.
Moreover, 90\% of the total value of exports was derived from agricultural produce \cite{Yearbook62}. 
This was a doubly perilous economic setting: one type of export and one major customer.
As early as 1961, Britain announced its interest in joining the European Economic Community (EEC),
and this was correctly interpreted in New Zealand as a serious warning that access to the UK market 
for agricutural products would sooner or later become difficult and that diversification was essential. By the time
Britain eventually joined the EEC in 1973, New Zealand had already become a completely independent
trading nation, as it remains today. In a remarkable change, by 1975 exports to the UK were down to 22\%
of the total and imports to 19\% \cite{Yearbook77}. 
As indicated above, IBM and other American computer companies took ample advantage of this opportunity.

\cbstart
Many New Zealanders experienced 1960-1975 as a prosperous era of growth,
despite a brief recession in 1967 \cite{Hall-2014} following a collapse in the export price of wool \cite{Easton-2020}.
Whereas the consumer price index rose by 124\% between 1960 and 1975, average wages grew by
204\%, according to the Reserve Bank of New Zealand’s inflation calculator.

Thus the growth of computing in New Zealand took place against the background of a
booming economy undergoing significant structural change, despite the Treasury’s continued
control over import licenses.
\cbend
The companies and other organizations that increasingly used computers during this period
presumably saw a strategic advantage in doing so. As noted above, project planning (PERT)
was a widespread application, which suggests that companies saw much more to computing than
savings on clerical staff. 

Writing of government computers in 1985, A. C. Shailes \cite{Shailes} said
“The use of computers has released staff from mundane administrative support tasks and
allowed them to become directly involved in achievement of the government’s main objectives.
It is clear that direct financial savings from computers have recovered the overall cost of all
computer resources many times over.” On the banking front, Ian Archibald \cite{Archibald} wrote
of Databank that “co-operation in basic processing systems has freed up staff to intensify competition
in developing and marketing the banking products.” Quantifying such effects in millions of dollars
seems impossible. Economists have largely failed to do so, even for intensively studied
economies such as the USA \cite{Oliner-1994}.
\cbstart
Nevertheless, we assert that the kind of competitive
and flexible economy that New Zealand needed to survive as a trading nation by 1975 would have been
much harder to achieve without information technology; the country could not have remained
competitive if it had failed to streamline its business practices in the same way as its trading partners.
\cbend
\vspace*{-8pt}
\section{SOCIAL IMPACT}
Perhaps it is no surprise that lawyers were among the first to consider the social impact of computers.
A first legal analysis was published in 1971 by Francis M. Auburn \cite{Auburn1971}, in which consideration of privacy
was the major theme. The author observed that “there is a fear of computers and their operators among the
general public” and quoted a rather naive counter-argument that “invasions of privacy are the result of misuses
of the data retrieved and not a function of the storage medium” (i.e., blame the users, not the technology).
With hundreds of thousands of banking transactions per day already taking place, and with a national police
computer and computerized hospital records in the offing, the risks of accidental or intentional privacy breaches
already seemed very real, more than 50 years ago. The legal situation was unclear: “If there is a general right of
privacy we are still no nearer to a definition of the circumstances in which it is applicable to computers.” All sorts
of data raised significant privacy issues once computerized, such as medical records, census returns, and credit history.
The issue of a unique national identity number to tie all records together was also of concern, with
inevitable references to George Orwell’s novel \textit{1984}. Auburn commented that “The most
rigorous presently known technological safeguards would be quite insufficient to
prevent abuse.” Auburn’s conclusion was that a “comprehensive review of the present and future
impact of electronic data processing on society” was needed, to be followed by legislation.

Auburn was not alone in raising these concerns. For example, also in 1971, A. L. C. Humphreys of ICL, the Christchurch
\textit{Press}, and the New Zealand Computer Society were all concerned about privacy. Certain lawyers
even drafted a proposed law \cite{Press-privacy-1971} and a Bill was introduced to Parliament in 1975
\cite{Press-privacy-1975} but never passed into law.
New Zealand’s first Privacy Act was not passed until 1993, prior to which common law applied.

Despite this legal vacuum, plans for a national police database were made from 1972, later enabled by Act of Parliament.
Because of ethical concerns,  the NZCS was highly involved in drafting the 1976 Wanganui Computer Centre Act.
The system, officially the National Law Enforcement Data Base, often called the Wanganui Computer after the town
where it was installed (now known as Whanganui), was operational by 1976 (Figure \ref{WanganuiCtr}).
This led some years later to protest actions. However, it seems that during our study period,
despite the concerns of some lawyers, there was little public perception of computers as a societal
threat prior to the police database. It was the latter which finally brought data privacy concerns to public attention.
\begin{figure}
\centerline{\includegraphics[width=18.5pc]{WanganuiCentre1975.jpg}}
\caption{\label{WanganuiCtr}Wanganui Computer Centre, 1975. Courtesy Fletcher Archives, NZ.}\vspace*{-5pt}
\end{figure}

Apart from privacy issues, it is hard to separate the impact of computers from other factors that
changed society in the 1960s and early 1970s, as the first generation born after World War II
grew to maturity. Colin Beardon \cite{Beardon} discussed  the computerization
of “women’s work” and its possible impact on unemployment. Certainly, skilled women who operated
ledger machines and comptometers found themselves out of a job as computers took over (Figure \ref{BNZ-ledger}).
Some of them no doubt moved into computer operations. However, many of the
replacement jobs would have been less skilled and lower paid data entry jobs.
Beardon noted that the fraction of the “unemployed who are female” increased from 10.9\% in 1961 to
18.0\% in 1971 and 34.9\% in 1975, but is impossible to know how much of this change was
caused by computers. For example, in retail banking, according to a retired banker \cite{Jonkers-private},
the technological change from manual to computerized operations changed people’s jobs,
but did not destroy them. The loss of typing and secretarial jobs came a decade
later, with the advent of personal computers. Throughout our period of study, computers were
mainly hidden from public view except for people directly involved, and had relatively little direct
social impact compared to other modern technology such as television. Apart from the job
market, where computers had a direct effect on some working lives, they were not even a
talking point for most people.
\begin{figure}
\centerline{\includegraphics[width=18.5pc]{BNZ-ledger-room-1950s.jpg}}
\caption{\label{BNZ-ledger}BNZ Ledger Room, 1950s. Courtesy BNZ Museum.}\vspace*{-5pt}
\end{figure}

Initial reactions to the introduction of computers were therefore based on curiosity. The first large
mainframes were used in only a few large organizations, and had little impact on the general public. 
During this period employment was high and although there was some discussion about whether 
computers would result in mass unemployment, generally they were viewed positively and there 
was optimistic discussion about shorter working weeks and an improved quality of life.
One concern was whether computers would improve individual autonomy or lead to increased 
conformity and control. Speaking in 1968, Databank’s General Manager, Gordon Hogg, felt that
automation would fit well with New Zealand values since a reduction in unskilled routine jobs would 
create a levelling in occupational social distinction, along with a continuing increase in income, both 
favored in a country committed to social equality \cite{Hogg-1968}. This reflected his experience at Databank where
 the new computer system did not lead to any job losses but rather released staff from routine 
functions like ledger posting and calculating interest so they were able to provide a more personal 
service to customers. 

The NZCS did take on the responsibility of considering the impact that computers would have on 
broader society. The third conference in 1972 had a theme of ‘Computers in the Community’ and
aimed to ensure that in New Zealand computers continued to be used for the greatest public benefit
and the least cost. The Computer Society’s president, Perce Harpham, remarked that this conference
“marks the maturing of our profession. We can no longer ignore the fact that our work affects the lives 
of many of our fellows and in turn that we are dependent on them” \cite{Harpham1972}. The fifth national conference in 
1976 also considered social impact; the theme was  ‘Responsibility’.  The keynote speaker, Paul Sieghart 
from the UK Data Protection Committee, stressed the responsibility that workers in the computing field 
had to educate the general public. Education is key, he said, as most people are ignorant about computing, 
regarding it “as a form of magic – if not sorcery”. The computer professional’s job was to demystify the 
computer in order to encourage the widest and most open public debate about all aspects of the social 
consequences of the developing technologies \cite{Sieghart1976}.
\cbstart
As noted above, this debate took place in the press as a direct result of NZCS activities.
\vspace*{-8pt}
\section{INTERNATIONAL COMPARISONS}
\cbend
New Zealand’s entry to the computer age was from some points of view not unique,
especially when considering other countries also in the process of cutting strong historical ties to the UK.
For example, the first computer in the Irish Republic was a BTM HEC-4 in 1957,
also known as the ICT 1201 \cite{Ireland-first}; such a machine was installed in
Wellington, New Zealand in 1960 \cite{FirstCinNZ}. Both Ireland and New Zealand
acquired an IBM 650 in 1960. The first Irish university computer was an IBM 1620, 
installed at University College Dublin in March 1962 \cite{Ireland-first}, two months before its twin
arrived at the University of Canterbury in New Zealand.

In India, the first computer imported was a HEC-2M, a direct ancestor of the HEC-4 or ICT 1201,
as early as 1955. The first machine in an Indian tertiary institution was an
IBM 1620 in 1963 \cite{Rajaraman-2015}. 

However, the Irish and Indian economies
both developed in completely different ways from that of New Zealand, despite
these superficial resemblances in the underlying equipment.
There was one similarity between the Indian and New Zealand
situations in the 1960s: in both countries, foreign exchange policy
severely limited computer imports. In contrast with New Zealand, both ICL 
and IBM therefore sought manufacturing licenses in India. IBM in particular successfully
manufactured equipment for the Indian market and for export \cite{Rajaraman-2015}.
\cbstart 
This was an opportunity that New Zealand might have had, but missed, probably due to its small domestic market.

\begin{table}[h]
\begin{center}
\begin{tabular}{ |p{0.13\linewidth}|p{0.1\linewidth}|p{0.1\linewidth}|p{0.1\linewidth}|p{0.1\linewidth}|p{0.1\linewidth}<{\raggedright}| } 
 \hline
Country & 1955 & 1960 & 1965 & 1970 & 1974\\
\hline
NZ & 0 & 0.8 & 22 & 57 & 93\\\hline
AU & 0.22 & 3.3 & 30 & 87 &173\\\hline 
UK & 0.25 & 4 & 29 & 113 & 256\\\hline
US & 1.4 & 30 & 127 & 361 & 772 \\\hline
FR & 0.1 & 3.6 & 30 & 105 & 300 \\\hline
DE (W) & 0.1 &5.4 &39 & 115 & 303\\\hline
JP & 0 & 0.9 & 19 & 85 & 237 \\\hline
\end{tabular}
\vspace{1ex}
\caption{\label{CpM}Computers per million residents.}
\end{center}
\end{table}

\begin{figure}
\centerline{\includegraphics[width=18.5pc]{ComputersPerHead.png}}
\caption{\label{CpH}Computers per million residents.}\vspace*{-5pt}
\end{figure}

Using data from Table~\ref{TotalComp}, \cite{Hendry-1989}, \cite{ACS-2017} and \cite{Australia-1974}, together with
generally available population data,  Table~\ref{CpM} and Figure~\ref{CpH} compare the dissemination of computers
relative to population size
in New Zealand with that in Australia, the UK, the USA, France, West Germany and Japan. Unsurprisingly,
the USA led the race. The European countries and Japan all had broadly similar profiles,
almost indistinguishable in the graph, but New Zealand
had fallen behind by 1974 with, for example, only 36\% as many computers per head as the UK.
In effect, New Zealand was about five years behind Europe in computerization
of the economy, and even further behind the USA.
It is also noticeable that Australia shows lower dissemination than the European group and Japan,
even though Australian scientists had built their own computers as early as the 1950s \cite{Ainsworth-2023}.
However, whatever New Zealanders might have felt at the time, objectively they
were falling behind the USA, Europe and Australia by the mid-1970s.

In the early 1970s, according to World Bank data, agriculture was about 13\% of New Zealand’s GDP, manufacturing
was about 25\%, and services 52\%. In contrast, the UK equivalents were agriculture 2\%, manufacturing 30\%,
and services 56\%. We hypothesize that the comparatively large GDP share of agriculture, where mainframe computing had
little role to play, allowed New Zealand to lag by a few years without serious consequences. Also, New Zealand lacked a defense industry of any size, a major driving factor in both the UK and the USA.

We can also ask whether New Zealand’s emphasis on service bureaus, caused by the Treasury’s import
license policy, was unusual.  One international study mentions that by 1965, the service sector was about 17\% of the
size of the computer systems market \cite{Campbell-Kelly-2008}. We do not have a corresponding data point for
New Zealand alone, but we know from \cite{ScottSurvey1967} that in 1967, there were about 100 computers
installed in total but at least 244 customers using bureau services. This suggests that the service bureau
model indeed represented significantly more than 17\% of the New Zealand computing market.

The first data privacy law in New Zealand was enacted only in 1993, nine years
after the UK’s Data Protection Act and 19 years after the first US Privacy Act. So here too,
New Zealand had nothing to boast about. In the academic and research domain, the country was
also about five years behind the USA and UK in creating Computer Science departments.

We have found no indication that the impact of computerization on the workforce was fundamentally
different in NZ compared to other countries. Contemporary studies are lacking, but for example
in 2021, 27\% of IT workers in NZ were female \cite{NZTIA-2022}, exactly the fraction observed
in government computing in 1974, and between the recent UK rate of 26\% \cite{WISE-2023}
and the Australian rate of 29\% \cite{Deloitte-2022}. It seems that after more than 60 years,
all three countries have reached a similar gender equilibrium in their IT workforces.

Finally, New Zealand’s examples of local ingenuity listed above are completely in accordance
with a national myth known as the “Number 8 Wire Mentality,” i.e. finding solutions to practical
problems using cheap and  locally available materials \cite{NZPOD, ingenuity}. However, there is nothing
exceptional about those examples, and as noted earlier, no local projects to build first-generation
computers are known, unlike Australia where there were at least two, CSIRAC and SILLIAC \cite{Philipson-2017}.

As already mentioned, New Zealand often claims that it is exceptional as a country, due to something more than its geographical remoteness. This claim particularly occurs in the popular press \cite{Frykberg-2025}. However, modern historians have taken a different view. For example, Fairburn \cite{Fairburn-2008} concludes that “what made New Zealand distinct is the
abnormal degree to which its people have borrowed from other cultures and
the particular combination of cultures they have borrowed from” (i.e., Australia, Britain and America).The editor of \textit{The New Oxford History of New Zealand} notes that its contributors “question the notion that New Zealand’s history is ‘unique’, distinct or exceptional.” \cite{Byrnes-2007} Apart from the lag in dissemination noted above, and the emphasis on service bureaus, we have indeed found nothing truly exceptional in New Zealand’s move into modern computing.
\cbend
\vspace*{-8pt}
\section{CONCLUSION}
We have described how computers penetrated government, business, tertiary education 
and research throughout New Zealand between 1960 and 1975. This was the period
when the country progressed from a post-colonial society strongly linked to the UK to
become a self-sufficient trading nation with its own identity. Modern computing,
local innovation and entrepreneurship combined to support the needs of the
evolving economy. Without this, it is unlikely that New Zealand could have
succeeded as it did. At the end of the period, computing was entering the public
consciousness as a technology that would soon have profound social implications,
although personal computing and social networking lay far in the future.
Developments since 1975 have been well documented by ITPNZ \cite{ITPNZ}
and by InternetNZ \cite{ConnectingClouds}.

\cbstart
However, we have found few indications that New Zealand was 
atypical among  a rather diverse group of countries. Computing started in the same
way in many countries. Indeed, New Zealand lagged a few years
behind its nearest neighbor (Australia), its former colonial power (the UK), other European
countries, and even behind Germany and Japan, countries recovering from massive wartime destruction.
New Zealand's claim to exceptionalism does not extend to the introduction of modern computing.
\cbend
\vspace*{-8pt}
\section{ACKNOWLEDGMENTS}
The National Library of New Zealand PapersPast archive
(\href{https://paperspast.natlib.govt.nz/newspapers}{https://paperspast.natlib.govt.nz/newspapers/})
and the very extensive Fletcher Trust Archives
(\href{https://collection.fletcherarchives.co.nz/}{https://collection.fletcherarchives.co.nz/}) were invaluable.
We are grateful to Robert MacCulloch, Professor of Economics at the University of Auckland, for his
insightful comments, 
\cbstart
and to Barbara Ainsworth, Monash University, for Australian data.
Informants with relevant
personal memories included Andy Bell, Tom Bell, Nevil Brownlee,
John Butcher, Neil Harsant, Hugh
Jonkers and Lloyd Thomas.
\cbend

%% style ieeetr doesn't display all details or alphabetise
%% style IEEEtran doesn't alphabetise
%% style plain doesn't meet IEEE requirements
\bibliographystyle{IEEEtranS}
\bibliography{growth}

\begin{IEEEbiography}{Brian E. Carpenter} is an Honorary Professor in the School of Computer Science, University of Auckland, New Zealand. He is interested in Internet protocol design and in computing history. He received a Ph.D. from the University of Manchester, UK and is a past Chair of the Internet Engineering Task Force. Contact: brian@cs.auckland.ac.nz\vspace*{8pt}
\end{IEEEbiography}

\begin{IEEEbiography}{Sathiamoorthy Manoharan} is a Senior Lecturer in the School of Computer Science, University of Auckland, New Zealand.
He is a senior member of IEEE. He is interested in computer systems, particularly with respect to performance and security. Contact: mano.manoharan@auckland.ac.nz.\vspace*{8pt}
\end{IEEEbiography}

\begin{IEEEbiography}{Janet Toland} is an Adjunct Professor in the School of Information Management,
Te Herenga Waka, Victoria University of Wellington, New Zealand. Her research work is in the history of
information systems, with a focus on social and ethical issues. Contact: Janet.Toland@vuw.ac.nz.\vspace*{8pt}
\end{IEEEbiography}



\end{document}

